\chapter{Introduction Générale}
\label{chap:intro}

\section{Contexte et Enjeux}

La cybersécurité moderne ne se limite plus à la technique. Elle est devenue un enjeu stratégique majeur pour la gouvernance des entreprises.

Pour les futurs ingénieurs et managers, trois compétences sont aujourd'hui indispensables :
\begin{enumerate}
    \item \textbf{Anticiper} : Identifier et évaluer les risques avant l'incident.
    \item \textbf{Vérifier} : Auditer les procédures et la conformité technique.
    \item \textbf{Réagir} : Gérer la crise et assurer la continuité de l'activité.
\end{enumerate}

Ce cours vous permettra de maîtriser le cycle de vie complet de la sécurité de l'information.

\section{Objectifs Pédagogiques}

À l'issue de ce cours, vous serez capable de :
\begin{itemize}
    \item Réaliser une analyse de risque complète (méthode EBIOS RM).
    \item Conduire un audit de sécurité (ISO 19011) et rédiger un rapport.
    \item Mettre en œuvre un plan de continuité d'activité (PCA).
    \item Piloter une cellule de crise et communiquer sous pression.
\end{itemize}

\section{Paysage Normatif}

Le cours s'appuie sur les standards internationaux reconnus. Ils constituent le langage commun des experts.

\begin{center}
\begin{tabular}{ll}
    \toprule
    \textbf{Domaine} & \textbf{Normes de référence} \\
    \midrule
    Gestion du Risque & \iso{31000}, \iso{27005}, EBIOS RM \\
    Management Sécurité & \iso{27001}, \iso{9001} \\
    Audit & \iso{19011}, COBIT \\
    Continuité & \iso{22301}, \iso{27035} \\
    \bottomrule
\end{tabular}
\end{center}

\section{Approche Pédagogique : Le Cycle PDCA}

Nous suivrons le cycle d'amélioration continue (PDCA - Plan Do Check Act) comme fil conducteur :

\begin{itemize}
    \item \textbf{Plan (Planifier) :} Analyse de risque et définition de la politique.
    \item \textbf{Do (Réaliser) :} Mise en œuvre des mesures techniques.
    \item \textbf{Check (Vérifier) :} Audit et mesure de l'efficacité.
    \item \textbf{Act (Améliorer) :} Gestion des incidents, crise et RETEX.
\end{itemize}

\section{Le Cas Pratique : Immersion}

La théorie ne suffit pas. Ce cours est articulé autour d'un cas d'entreprise unique et réaliste : \textbf{LiZbiZ}.

\begin{boxlizbiz}
Vous ne serez pas de simples spectateurs. Vous agirez en tant que consultants ou membres de la DSI de l'entreprise \textbf{LiZbiZ}. 
    
Un laboratoire virtuel (GNS3) simulant leur Système d'Information vous sera fourni. Chaque module du cours s'appuiera sur ce cas concret pour appliquer les concepts théoriques.
\end{boxlizbiz}

Le chapitre suivant présente en détail l'entreprise LiZbiZ et son architecture technique.


% ==============================================================================
% Fichier : 01_cas_lizbiz.tex
% Description : Présentation complète du cas pratique fil rouge
% ==============================================================================
